% Options for packages loaded elsewhere
\PassOptionsToPackage{unicode}{hyperref}
\PassOptionsToPackage{hyphens}{url}
\documentclass[
]{article}
\usepackage{xcolor}
\usepackage{amsmath,amssymb}
\setcounter{secnumdepth}{-\maxdimen} % remove section numbering
\usepackage{iftex}
\ifPDFTeX
  \usepackage[T1]{fontenc}
  \usepackage[utf8]{inputenc}
  \usepackage{textcomp} % provide euro and other symbols
\else % if luatex or xetex
  \usepackage{unicode-math} % this also loads fontspec
  \defaultfontfeatures{Scale=MatchLowercase}
  \defaultfontfeatures[\rmfamily]{Ligatures=TeX,Scale=1}
\fi
\usepackage{lmodern}
\ifPDFTeX\else
  % xetex/luatex font selection
\fi
% Use upquote if available, for straight quotes in verbatim environments
\IfFileExists{upquote.sty}{\usepackage{upquote}}{}
\IfFileExists{microtype.sty}{% use microtype if available
  \usepackage[]{microtype}
  \UseMicrotypeSet[protrusion]{basicmath} % disable protrusion for tt fonts
}{}
\makeatletter
\@ifundefined{KOMAClassName}{% if non-KOMA class
  \IfFileExists{parskip.sty}{%
    \usepackage{parskip}
  }{% else
    \setlength{\parindent}{0pt}
    \setlength{\parskip}{6pt plus 2pt minus 1pt}}
}{% if KOMA class
  \KOMAoptions{parskip=half}}
\makeatother
\usepackage{color}
\usepackage{fancyvrb}
\newcommand{\VerbBar}{|}
\newcommand{\VERB}{\Verb[commandchars=\\\{\}]}
\DefineVerbatimEnvironment{Highlighting}{Verbatim}{commandchars=\\\{\}}
% Add ',fontsize=\small' for more characters per line
\newenvironment{Shaded}{}{}
\newcommand{\AlertTok}[1]{\textcolor[rgb]{1.00,0.00,0.00}{\textbf{#1}}}
\newcommand{\AnnotationTok}[1]{\textcolor[rgb]{0.38,0.63,0.69}{\textbf{\textit{#1}}}}
\newcommand{\AttributeTok}[1]{\textcolor[rgb]{0.49,0.56,0.16}{#1}}
\newcommand{\BaseNTok}[1]{\textcolor[rgb]{0.25,0.63,0.44}{#1}}
\newcommand{\BuiltInTok}[1]{\textcolor[rgb]{0.00,0.50,0.00}{#1}}
\newcommand{\CharTok}[1]{\textcolor[rgb]{0.25,0.44,0.63}{#1}}
\newcommand{\CommentTok}[1]{\textcolor[rgb]{0.38,0.63,0.69}{\textit{#1}}}
\newcommand{\CommentVarTok}[1]{\textcolor[rgb]{0.38,0.63,0.69}{\textbf{\textit{#1}}}}
\newcommand{\ConstantTok}[1]{\textcolor[rgb]{0.53,0.00,0.00}{#1}}
\newcommand{\ControlFlowTok}[1]{\textcolor[rgb]{0.00,0.44,0.13}{\textbf{#1}}}
\newcommand{\DataTypeTok}[1]{\textcolor[rgb]{0.56,0.13,0.00}{#1}}
\newcommand{\DecValTok}[1]{\textcolor[rgb]{0.25,0.63,0.44}{#1}}
\newcommand{\DocumentationTok}[1]{\textcolor[rgb]{0.73,0.13,0.13}{\textit{#1}}}
\newcommand{\ErrorTok}[1]{\textcolor[rgb]{1.00,0.00,0.00}{\textbf{#1}}}
\newcommand{\ExtensionTok}[1]{#1}
\newcommand{\FloatTok}[1]{\textcolor[rgb]{0.25,0.63,0.44}{#1}}
\newcommand{\FunctionTok}[1]{\textcolor[rgb]{0.02,0.16,0.49}{#1}}
\newcommand{\ImportTok}[1]{\textcolor[rgb]{0.00,0.50,0.00}{\textbf{#1}}}
\newcommand{\InformationTok}[1]{\textcolor[rgb]{0.38,0.63,0.69}{\textbf{\textit{#1}}}}
\newcommand{\KeywordTok}[1]{\textcolor[rgb]{0.00,0.44,0.13}{\textbf{#1}}}
\newcommand{\NormalTok}[1]{#1}
\newcommand{\OperatorTok}[1]{\textcolor[rgb]{0.40,0.40,0.40}{#1}}
\newcommand{\OtherTok}[1]{\textcolor[rgb]{0.00,0.44,0.13}{#1}}
\newcommand{\PreprocessorTok}[1]{\textcolor[rgb]{0.74,0.48,0.00}{#1}}
\newcommand{\RegionMarkerTok}[1]{#1}
\newcommand{\SpecialCharTok}[1]{\textcolor[rgb]{0.25,0.44,0.63}{#1}}
\newcommand{\SpecialStringTok}[1]{\textcolor[rgb]{0.73,0.40,0.53}{#1}}
\newcommand{\StringTok}[1]{\textcolor[rgb]{0.25,0.44,0.63}{#1}}
\newcommand{\VariableTok}[1]{\textcolor[rgb]{0.10,0.09,0.49}{#1}}
\newcommand{\VerbatimStringTok}[1]{\textcolor[rgb]{0.25,0.44,0.63}{#1}}
\newcommand{\WarningTok}[1]{\textcolor[rgb]{0.38,0.63,0.69}{\textbf{\textit{#1}}}}
\usepackage{longtable,booktabs,array}
\newcounter{none} % for unnumbered tables
\usepackage{calc} % for calculating minipage widths
% Correct order of tables after \paragraph or \subparagraph
\usepackage{etoolbox}
\makeatletter
\patchcmd\longtable{\par}{\if@noskipsec\mbox{}\fi\par}{}{}
\makeatother
% Allow footnotes in longtable head/foot
\IfFileExists{footnotehyper.sty}{\usepackage{footnotehyper}}{\usepackage{footnote}}
\makesavenoteenv{longtable}
\setlength{\emergencystretch}{3em} % prevent overfull lines
\providecommand{\tightlist}{%
  \setlength{\itemsep}{0pt}\setlength{\parskip}{0pt}}
\usepackage{hyperxmp}
\usepackage{bookmark}
\IfFileExists{xurl.sty}{\usepackage{xurl}}{} % add URL line breaks if available
\urlstyle{same}
\hypersetup{
  pdftitle={Preserving semantic continuity across actors: a tell-based approach without orchestration},
  pdfauthor={Alvaro Rivera},
  pdfkeywords={\xmpquote{actor model}, \xmpquote{semantic
continuity}, \xmpquote{cross-actor causation}, \xmpquote{message
passing}, \xmpquote{design theory}, \xmpquote{journal-based
program}, \xmpquote{actor systems}, \xmpquote{tell
primitive}, \xmpquote{puppeteer framework}},
  hidelinks,
  pdfcreator={LaTeX via pandoc}}

\title{Preserving semantic continuity across actors: a tell-based
approach without orchestration}
\author{Alvaro Rivera}
\date{2026-05-10}

\begin{document}
\maketitle
\begin{abstract}
This paper is a design theory contribution in the sense of Alan Hevner,
Stuart March, Jinsoo Park, and Sudha Ram (2004): it identifies a
structural defect in the canonical actor-systems literature, derives the
design principles required to address it, and presents an instantiation
demonstrating that the alternative is realizable in production. For five
decades, from Carl Hewitt's original formulation through Gul Agha, Joe
Armstrong, and modern frameworks such as Akka and Microsoft Orleans,
cross-actor causation has been treated as an operational concern of the
runtime rather than as a construct of the program. As a result, the
program of an actor is structurally fragmented at every actor boundary:
when one actor causes an effect in another, the causal sentence
disappears from the program and reappears as infrastructure. This paper
shows that this fragmentation is not inherent to the actor model but a
contingent historical choice. It derives the design principles under
which cross-actor message passing can be expressed as a sentence of the
actor's program without violating actor isolation and without
introducing orchestration. An existing primitive, tell, is shown to
satisfy these principles when recorded as part of the actor's journal.
Under this construction, program continuity extends across actors while
preserving the defining properties of actor systems.
\end{abstract}

\section{Preserving semantic continuity across
actors}\label{preserving-semantic-continuity-across-actors}

\subsection{TL;DR}\label{tldr}

\begin{quote}
The actor model treats cross-actor causation as operational --- message
dispatch by a runtime, not a statement in any actor's program. The
treatment is not entailed by the model; it is a contingent design
decision adopted in the canonical sources fifty years ago and never
seriously challenged. The cost is structural: an entire ecosystem of
compensating patterns --- sagas, choreography, distributed tracing,
workflow engines --- exists to reconstruct cross-actor causal chains
that no actor's program records.

This paper observes that the assumption can be removed without altering
any structural commitment of the actor model. Three conditions describe
a system in which the cross-actor send is recorded as a sentence in the
sender's program: locality of writes (each journal records only its own
actor's activity), causation as program statement (the cross-actor send
is a sentence in the sender's journal), and no external coordinator (no
party outside the participating actors decides what happens next). Any
primitive satisfying these three conditions dissolves the assumption.
\emph{Tell} --- a primitive in the Puppeteer framework --- is one such
realization.

Three labs are exhibited side by side on the same domain (saga,
choreography, tell). The reader sees, in journals, that the saga places
the joint history in a coordinator's program, that the choreography
places it in an external bus log, and that the tell instantiation places
it in the sender's own program. Three additional tests demonstrate that
under tell the journal alone supports replay reconstruction,
cross-datacenter replication, and audit query --- operations that
require external apparatus under the assumption. \textbf{Forty years of
convention are not forty years of necessity.}
\end{quote}

\begin{center}\rule{0.5\linewidth}{0.5pt}\end{center}

\subsection{Claims this paper makes}\label{claims-this-paper-makes}

\begin{enumerate}
\def\labelenumi{\arabic{enumi}.}
\item
  \textbf{The assumption named.} Cross-actor causation has been treated
  as operational rather than programmatic across fifty years of
  canonical actor-systems literature --- from Hewitt 1973 through Akka
  and Orleans. \emph{(Verification: §2 + the genealogy table in §2.4.)}
\item
  \textbf{The assumption is contingent.} No theorem of the actor model
  entails it. The three structural commitments of the model ---
  autonomy, message-based communication, isolation --- define the
  ontology of actors, not the mechanics of their runtimes.
  \emph{(Verification: §4.)}
\item
  \textbf{Three conditions resolve the assumption.} Locality of writes
  (C1), causation as program statement (C2), and no external coordinator
  (C3) describe a system in which semantic continuity is preserved
  across actors without violating any of the actor model's commitments.
  The conditions are not design choices; they are the minimal
  consequences of reconciling the model's commitments with the
  construct. \emph{(Verification: §7.1.)}
\item
  \textbf{Any primitive satisfying C1+C2+C3 dissolves the assumption.}
  The contribution is conceptual; the realization is one of many
  possible. An actor framework whose programs do not currently record
  cross-actor sends could, in principle, be extended to satisfy the
  three conditions. \emph{(Verification: §7.5.)}
\item
  \textbf{The actor model's commitments survive the reformulation
  intact.} Autonomy, message-based communication, and isolation each
  remain unchanged; only the historical interpretation of what each
  actor's program is \emph{permitted to say} is removed.
  \emph{(Verification: §7.3.)}
\item
  \textbf{The compensating ecosystem exists because of the assumption.}
  Saga orchestrators, event-driven choreography, distributed tracing,
  and workflow engines each compensate, in their own way, for the
  absence of cross-actor causation in any actor's program. Their
  sophistication, maturity, and widespread adoption are evidence of the
  cost of the assumption, not of its correctness. \emph{(Verification:
  §6 + the side-by-side labs in §8.3.)}
\item
  \textbf{Tell exhibits the conditions in production.} A Reaction whose
  \texttt{.Causation.Continue(...)} body issues a \texttt{tell} writes
  one journal entry on the sender's side; the receiver's acknowledgment
  closes the round-trip in a second entry. The sender's journal becomes
  a self-contained record of what crossed the actor boundary.
  \emph{(Verification: §8.2 + §8.3 Style 3.)}
\item
  \textbf{Auditing the cross-actor narrative becomes reading the
  program.} The journal shows what was sent, to whom, with what content,
  at what time, and with what acknowledgment --- without correlation
  IDs, distributed tracing, or external aggregation.
  \emph{(Verification: §8.5 G3.)}
\item
  \textbf{Replay reconstructs cross-actor state from the journal alone.}
  A fresh actor with no shared transport and no live receiver, replaying
  the journal, reconstructs the in-flight tell state.
  \emph{(Verification: §8.5 G1.)}
\item
  \textbf{Cross-datacenter replication preserves the cross-actor causal
  chain.} The journal carries the cross-actor causation across data
  centers because the causation was always recorded in a place
  replication can carry. \emph{(Verification: §8.5 G2.)}
\end{enumerate}

\begin{center}\rule{0.5\linewidth}{0.5pt}\end{center}

\subsection{1. Introduction}\label{introduction}

This paper makes a design theory contribution. It identifies and names a
structural assumption that the canonical literature on actor-based
systems has repeatedly documented in different forms without recognizing
as a single construct; derives the design principles required to reject
it; and presents an instantiation --- a system in which those principles
have been realized in production --- as confirmation that the construct
is realizable. The contribution is conceptual; the instantiation serves
as an existence proof, not as the substance of the claim. The genre is
that described by Alan Hevner, Stuart March, Jinsoo Park, and Sudha Ram
(2004) as design science research: empirical evidence is provided in the
form of a working artifact rather than a controlled experiment.

This is not a systems paper. It does not present performance benchmarks,
fault-injection metrics, or latency comparisons against existing actor
frameworks. The genre --- design theory --- measures contribution by the
precision of the construct, the validity of the principles, and the
realizability of the instantiation. Section 8 exhibits the instantiation
and tests whether the mechanism satisfies its stated structural
properties. Readers expecting quantitative comparisons against
alternatives will find structural comparisons --- what each pattern
records, where the joint history lives --- in §6 and §8.4.

Consider an ordinary observation about actor-based systems. An actor
performs an operation and, as a result, sends a message to another
actor. The first actor's operation, state mutation, and emitted events
appear in its journal or trace. The second actor receives the message
and its own operation is likewise recorded. Yet the act of sending the
message --- the causal step that connects the two --- does not appear in
either actor's program. It is mediated by the runtime, the dispatcher,
or the message broker. The sender's journal does not record that it
spoke; the receiver's journal does not record, in program terms, who
spoke to it.

This observation is so familiar that it rarely appears noteworthy. But
it should. If actors are programs, and a sentence in one actor's program
causes an effect in another, then that causal sentence has every reason
to appear as part of the first actor's program. That it does not is not
a logical consequence of the actor model; it is a structural feature of
how actor systems have historically been constructed.

This paper names that gap. The construct introduced is \emph{semantic
continuity}: the property of a program (in the substrate-level sense
established in Paper 2 §1.2: the pair of domain library and journal of
invocations) whose causal structure remains recorded as part of the
program itself, even when its effects cross boundaries. The defect
identified is the absence of semantic continuity at actor boundaries.
The principles derived are the conditions under which semantic
continuity can be preserved without violating actor isolation and
without introducing orchestration. The instantiation presented is
\emph{tell}, a primitive in which the cross-actor send is recorded as a
sentence of the sender's journal.

§2 traces the genealogy of the assumption that produces this gap,
showing that across five decades and multiple canonical generations of
literature the assumption has remained continuous in substance though
varied in form. §3 names the assumption explicitly: \emph{causation
between actors is treated as operational rather than programmatic}. §4
demonstrates that this assumption is contingent rather than necessary
--- no theorem of the actor model entails it. §5 examines the
architectural and operational consequences of the assumption. §6 shows
why existing responses to those consequences --- sagas, choreography,
distributed tracing, and workflow engines --- cannot dissolve the
assumption, because they reconstruct cross-actor flow after the fact
rather than preserving it as program. §7 reformulates the model: under
three explicit conditions, semantic continuity can be preserved across
actors. §8 presents the instantiation, \emph{tell}, through a
comparative case study against saga and choreography implementations of
the same domain. §9 relates the construct to prior work in this paper
series. §10 concludes.

\begin{center}\rule{0.5\linewidth}{0.5pt}\end{center}

\subsection{2. Genealogy}\label{genealogy}

\subsubsection{2.1 The founding decision
(1970s)}\label{the-founding-decision-1970s}

The actor model was introduced in 1973 by Hewitt, Bishop, and Steiger as
a unifying account of concurrent computation in which ``all of the modes
of behavior can be defined in terms of one kind of behavior: sending
messages to actors'' {[}Hewitt et al.~1973, p.~235{]}. Sending messages
is positioned as a universal primitive, but the framing goes further: it
is positioned as a \emph{machine-level} primitive. ``The basic unit of
execution on an actor machine is sending a message in much the same way
that the basic unit of execution on present day machines is an
instruction'' {[}ibid., p.~236{]}. The act of sending is therefore
parallel to a hardware instruction --- it is what the abstract machine
performs \emph{between} actors, not content of any actor's program.

A parallel tradition emerged in 1978 with Hoare's \emph{Communicating
Sequential Processes} {[}Hoare 1978{]}. Where the actor model offered
messages between named actors, CSP offered input and output as basic
primitives of programming and synchronized parallel composition as a
fundamental program structuring method. Processes are independent; their
communication occurs at named ports through a symmetric handshake. The
handshake is a feature of the language, not a statement in either
process's program.

The actor model was formalized in 1986 in Agha's MIT thesis {[}Agha
1986{]}, which defined actors as behavior functions over messages and
treated cross-actor effects as emitted output messages of the function.
The formal account ratified what Hewitt 1973 had stated informally: an
actor's program describes how it responds to the messages it receives;
what happens \emph{between} actors is the operational semantics of the
model, not the content of any actor's program.

By the close of the 1980s, two foundational traditions --- actors and
CSP --- had stabilized the same architectural decision under different
vocabularies. The boundary between an actor (or process) and the
message-passing layer had been drawn deliberately, and the layer between
actors had been characterized as operational rather than programmatic.

\subsubsection{2.2 The pragmatic maturation
(1990s--2000s)}\label{the-pragmatic-maturation-1990s2000s}

The operational frame became productive in Erlang. Armstrong's 2003 PhD
thesis {[}Armstrong 2003{]} argued that fault tolerance becomes
tractable precisely because processes do not share programmatic flow:
failures are handled by other processes via supervision links, ``ways
for one process to react to the failure of another, supported directly
by the VM'' {[}Wensel 2018, reflecting on Armstrong's
architectural-model chapters{]}. Supervision is a structural pattern in
which a supervisor observes the failure of a child process and reacts.
The supervisor's program is local; the failed process's program does not
extend into the supervisor; the cross-process relationship is mediated
by VM-level links and monitors. Cross-process causation is an
infrastructure concern. The frame stabilized further by becoming useful:
the operational-rather-than-programmatic separation was now what made
fault-tolerant systems possible.

\subsubsection{2.3 The modern instantiations
(2010s--)}\label{the-modern-instantiations-2010s}

Two implementations carried the actor frame into modern production
systems. Akka, the canonical JVM actor framework {[}Lightbend,
\emph{Akka core: Interaction Patterns}{]}, characterizes its fundamental
message-send pattern directly:

\begin{quote}
\emph{``Tell is asynchronous which means that the method returns right
away. After the statement is executed there is no guarantee that the
message has been processed by the recipient yet. It also means there is
no way to know if the message was received, the processing succeeded or
failed.''}
\end{quote}

The verb is named --- \texttt{tell} --- and so is its absence of
program-level effect. The dispatch occurs; the sender's program does not
record that it occurred, nor whether the recipient acted on it.

Microsoft's Orleans {[}Bernstein et al.~2014{]} introduced \emph{virtual
actors}: location-transparent grains whose dispatch is mediated by the
Orleans runtime. Each grain has a key; the runtime decides where the
grain lives and routes calls; the grain's code observes only local state
and method invocation. The runtime maintains the cross-grain dispatch
logic as infrastructure; the grain's program never sees nor records
cross-grain causation as part of its narrative.

\begin{quote}
\textbf{Forty years apart, one assumption.} In Hewitt 1973:
\emph{``Sending a message to an actor makes no presupposition that the
actor sent the message will ever send back a message to the
continuation''} (p.~241). In Akka 2014--present: \emph{``there is no way
to know if the message was received, the processing succeeded or
failed.''} The two formulations are forty years apart and identical in
substance. The verb \texttt{tell} has been part of actor-systems
vocabulary throughout. What this paper observes is that the verb names
the dispatch, not the program: a \texttt{tell} from actor A to actor B
leaves no trace in either actor's program. The verb has existed; the
assumption that the verb's effect was extra-programmatic has survived
intact alongside it.
\end{quote}

\subsubsection{2.4 The naturalized
assumption}\label{the-naturalized-assumption}

Across fifty years and four canonical generations of literature ---
Hewitt's foundational paper, Hoare's parallel tradition, Agha's
formalization, Armstrong's pragmatic maturation, and the modern
instantiations in Akka and Orleans --- the question \emph{``Is
cross-actor flow part of any actor's program?''} is not asked. It does
not need to be asked, because the answer has been assumed. The
assumption is so stabilized that it does not appear as a defended
thesis; it appears as the default frame within which every subsequent
contribution operates.

The forms in which the assumption surfaces vary by generation; the
substance is continuous.

{\def\LTcaptype{none} % do not increment counter
\begin{longtable}[]{@{}
  >{\raggedright\arraybackslash}p{(\linewidth - 6\tabcolsep) * \real{0.2500}}
  >{\raggedright\arraybackslash}p{(\linewidth - 6\tabcolsep) * \real{0.2500}}
  >{\raggedright\arraybackslash}p{(\linewidth - 6\tabcolsep) * \real{0.2500}}
  >{\raggedright\arraybackslash}p{(\linewidth - 6\tabcolsep) * \real{0.2500}}@{}}
\toprule\noalign{}
\begin{minipage}[b]{\linewidth}\raggedright
Year
\end{minipage} & \begin{minipage}[b]{\linewidth}\raggedright
Work
\end{minipage} & \begin{minipage}[b]{\linewidth}\raggedright
Anchor
\end{minipage} & \begin{minipage}[b]{\linewidth}\raggedright
Form of the assumption
\end{minipage} \\
\midrule\noalign{}
\endhead
\bottomrule\noalign{}
\endlastfoot
1973 & Hewitt et al.~--- \emph{A Universal Modular ACTOR Formalism}
(IJCAI) & ``The basic unit of execution on an actor machine is sending a
message in much the same way that the basic unit of execution on present
day machines is an instruction.'' (p.~236) & Sending is a machine-level
primitive, parallel to a hardware instruction; therefore not content of
any actor's program. \\
1978 & Hoare --- \emph{Communicating Sequential Processes} (CACM) &
Input and output are basic primitives of programming; parallel
composition is a fundamental program structuring method. &
Process-to-process synchronization is a feature of the language's
parallel composition, not a statement in either process's program. \\
1986 & Agha --- \emph{Actors: A Model of Concurrent Computation} (MIT) &
Actors are behavior functions; cross-actor effects are emitted output
messages of the function. & The behavior function is local; the
cross-actor effect is the \emph{output} of the function, separate from
the function's body. \\
2003 & Armstrong --- \emph{Making reliable distributed systems\ldots{}}
(KTH) & ``All failures are handled via `monitors' and `links', which are
ways for one process to react to the failure of another, supported
directly by the VM.'' & Cross-process flow is infrastructure-level (VM
links/monitors), not part of any process's program. \\
2010s & Lightbend --- Akka \emph{Interaction Patterns} & ``There is no
way to know if the message was received, the processing succeeded or
failed.'' & The dispatch verb (\texttt{tell}) exists; its effect is
explicitly outside the sender's program --- no acknowledgment record. \\
2014 & Bernstein et al.~--- \emph{Orleans: Distributed Virtual Actors} &
The Orleans runtime places grains and mediates cross-grain calls; grain
code sees only local state. & Cross-grain dispatch is hidden by the
runtime; the grain's program never sees nor records cross-grain
causation. \\
\end{longtable}
}

The continuity of this assumption across five decades is not the result
of oversight, but of success: actor systems worked precisely because the
separation between program and message-passing layer was operationally
effective. What follows does not dispute that effectiveness; it
questions whether the separation is necessary.

§3 names this assumption explicitly. §4 demonstrates that it is
contingent rather than necessary.

\begin{center}\rule{0.5\linewidth}{0.5pt}\end{center}

\subsection{3. The assumption named}\label{the-assumption-named}

What §2 has shown to be continuous across fifty years admits a single
one-line formulation. The structural claim that links all six entries of
§2.4 --- Hewitt's ``machine-level instruction'', Hoare's parallel
composition, Agha's emitted output messages, Armstrong's VM-supported
links, Akka's \texttt{tell}, Orleans' runtime-mediated dispatch --- is:

\begin{quote}
\textbf{Causation between actors is operational, not programmatic.}
\end{quote}

The two terms in the formulation are the working categories of the rest
of this paper.

\emph{Operational} designates effects that the system performs around or
between actors but that no actor's program records. Message dispatch by
a runtime, supervision links between processes, virtual-actor placement
decisions, message-broker routing, and distributed-tracing correlation
IDs are all operational in this sense: their existence is mediated by
infrastructure, and their effect on the system is real, but no program
in the system contains the act of mediation as a statement.

\emph{Programmatic} designates effects that an actor's program contains
as a statement --- a line of code, a journal entry, a recorded behavior
--- in such a way that the effect is observable from within the
program's own narrative. A method invocation that an actor performs on
itself is programmatic in this sense; so is a state mutation written by
that actor; so is a journal entry recorded by that actor. The actor's
program records what happened; replay of that program reconstructs what
happened.

The contrast can be visualized:

\begin{verbatim}
Under the assumption (§3): causation operational, not programmatic

    A's program                                     B's program
        │                                                 ▲
        │      runtime / broker / dispatcher              │
        │  ─────────────────────────────────────────────► │
        │      (cross-actor causation not                 │
        │       recorded as program)                      │
        ▼                                                 ▼
    A's journal                                     B's journal
    (local effects only)                            (local effects only)


Under the alternative (§7): causation as program statement

    A's program                                     B's program
        │                                                 ▲
        │  ─────────────────  tell  ────────────────────► │
        │                                                 │
        ▼                                                 ▼
    A's journal                                     B's journal
    [tell entry]                                    [receipt entry]
    [ack entry]
\end{verbatim}

Both views describe the same flow; they differ in what is recorded as
program. Under the assumption, the cross-actor send lives in
infrastructure between actors and is invisible from any actor's program.
Under the alternative, the cross-actor send is itself a journal entry on
the sender's side, with a corresponding receipt entry on the receiver's
side.

The naturalized assumption is the claim that, in the ontology of an
actor system, the act of one actor causing an effect in another belongs
to the \emph{operational} category and not the \emph{programmatic} one.
The 1973 framing of message-sending as a hardware-instruction analogue
is this claim. The 2014--present framing of \texttt{tell} as a
fire-and-forget dispatch with no acknowledgment record is this claim.
Across formulations, the assumption holds.

§4 demonstrates that the assumption, while continuous in the literature,
is contingent rather than necessary: no theorem of the actor model
entails it. §5 examines the consequences of the assumption --- what is
structurally lost when cross-actor causation lives outside any program.
§6 shows why the existing repertoire of patterns built atop the actor
model --- sagas, choreography, distributed tracing, workflow engines ---
cannot dissolve the assumption, because they reconstruct cross-actor
flow \emph{after} the fact rather than preserving it as program. §7
reformulates: the assumption is contingent, the alternative is
constructible, and the cross-actor flow can be a sentence of the
program.

\begin{center}\rule{0.5\linewidth}{0.5pt}\end{center}

\subsection{4. Contingency}\label{contingency}

The assumption named in §3 --- that causation between actors is treated
as operational rather than programmatic --- would be necessary if it
followed from a theorem of the actor model, from a foundational property
the model formally entails, or from the structural commitments that
distinguish actor systems from other concurrency paradigms. It does not.
This section traces what the actor model formally requires and what it
does not, and shows that the assumption is a contingent feature of how
actor systems have historically been built rather than a consequence of
what they are.

\subsubsection{4.1 What the actor model
entails}\label{what-the-actor-model-entails}

The actor model, in its canonical formulations from Hewitt 1973 through
Agha 1986, formally entails three structural commitments:

\begin{enumerate}
\def\labelenumi{\arabic{enumi}.}
\tightlist
\item
  \textbf{Autonomy.} Each actor has its own state and processes messages
  serially. No actor's behavior is contingent on the simultaneous
  execution of another's.
\item
  \textbf{Message-based communication.} Actors communicate by passing
  messages. No actor reads or writes another's state directly.
\item
  \textbf{Isolation.} An actor's state is private to it. The address
  space of an actor is not shared.
\end{enumerate}

These three commitments are what distinguish actor systems from
shared-memory concurrency models. They define the ontology of actors,
not the mechanics of their runtimes. They are the substance of ``what
the actor model is'' in any reading of the canonical sources.

\subsubsection{4.2 What it does not
entail}\label{what-it-does-not-entail}

The three commitments above do not entail that the \emph{act of sending}
be invisible to the sender's own program. They do not entail that
cross-actor dispatch be mediated by a runtime layer that owns the send.
They do not entail that a record of ``actor A spoke to actor B at time
T'' must reside outside actor A's narrative.

In Hewitt 1973, the act of sending is an action the actor performs. The
framing of message-sending as a machine-level instruction analogous to a
hardware instruction {[}Hewitt et al.~1973, p.~236{]} is an interpretive
choice that supports the architecture proposed in that paper; it is not
a formal consequence of what an actor is. In Agha 1986, the actor's
behavior is modelled as a function from (state, message) to (next state,
outgoing messages) {[}Agha 1986, ch.~4{]}. The outgoing messages are
output of the function. The formalization neither requires nor prohibits
the function from producing, in addition, a record of the send
observable in the actor's own program --- the question is not raised.
The formalization can be extended, without modifying any of the three
commitments above, to include as output not only the outgoing messages
but also a record of the send observable within the actor's own program.

What the canonical formulations establish is the \emph{boundary} between
actors --- autonomy, isolation, message passing as the exclusive
cross-actor mechanism. They do not establish that the boundary be
invisible from within. The opacity of the cross-actor send to the
sender's program is a separate decision that the canonical sources
adopted but did not derive.

\subsubsection{4.3 The objection from
isolation}\label{the-objection-from-isolation}

The most common objection to making cross-actor sends programmatic is
that doing so would violate actor isolation. The objection conflates two
distinct properties.

Isolation requires that no actor read or write another's state.
Recording a cross-actor send in the sender's journal does not require
either. The sender writes to its own journal --- its own state. The
journal entry describes what the sender did: \emph{``I sent a message of
type M to actor B at time T.''} The receiver, when it processes the
message, writes a corresponding entry to its own journal: \emph{``I
received a message of type M from actor A at time T.''} Both records are
local; neither requires shared state. Isolation is preserved.

What changes is not who can read whose state, but what each actor's
program is permitted to say about its own activity. That change is local
to the actor's own program; it is invisible across the boundary.

\subsubsection{4.4 The objection from
orchestration}\label{the-objection-from-orchestration}

The second common objection is that recording cross-actor causation as
program would introduce orchestration. The objection conflates a record
of causation with a director of behavior.

Orchestration is the architectural pattern in which an external
coordinator decides what each participating actor does next. The
coordinator dispatches commands; the participants execute them.
Causation flows from the coordinator outward.

A journal entry that records what an actor did is not a coordinator. It
does not direct anything. Each actor still processes its own messages
autonomously, decides what to do, and emits its own outgoing messages.
The journal preserves what happened; an orchestrator would prescribe
what should happen. Recording causation does not introduce a director.
It introduces narration, not control.

\subsubsection{4.5 Closing}\label{closing}

The assumption stated in §3 is therefore not entailed by the actor
model. It does not follow from Hewitt's formulation, from Agha's
formalization, from the requirement of isolation, or from the absence of
orchestration. The opacity of the cross-actor send to the sender's
program is a contingent design decision --- adopted by the canonical
sources, propagated through the lineage in §2, and productive enough to
remain unexamined. It is not, in any formal sense, a property the model
requires.

The consequence is decisive: the absence of semantic continuity at actor
boundaries is not a necessity of the actor model, but a historical
artifact of how actor systems have been constructed.

\begin{center}\rule{0.5\linewidth}{0.5pt}\end{center}

\subsection{5. Consequences}\label{consequences}

When the assumption stated in §3 is in force, cross-actor causation is
held outside any actor's program. The consequences of this displacement
are not operational inconveniences; they are structural. Four are
examined here. Each is a distinct symptom; each is the same absence ---
the absence of semantic continuity at actor boundaries --- observed from
a different angle.

\subsubsection{5.1 Auditability requires external
reconstruction}\label{auditability-requires-external-reconstruction}

The question \emph{``why did this happen?''}, asked of an effect that
crossed an actor boundary, cannot be answered by reading any single
actor's program. The sender's journal records what the sender did; the
receiver's journal records what the receiver did; neither records the
link. To answer the question, an auditor must consult tools that live
outside the programs: correlation IDs threaded through messages by the
runtime, distributed tracing systems that observe the runtime from
outside, log aggregation pipelines that join records across actor
boundaries.

These tools work, but their necessity is evidence that the program does
not contain what they reconstruct. They do not extend the program; they
substitute for it. Auditability therefore depends on infrastructure that
compensates for what the program does not record.

\subsubsection{5.2 Replay does not reconstruct cross-actor
history}\label{replay-does-not-reconstruct-cross-actor-history}

Replaying an actor's journal reconstructs that actor's local history:
which messages it received, how it responded, what state it produced.
Replay of two actors' journals reconstructs two local histories. It does
not reconstruct the joint history. The joint history never existed as a
program artifact to be replayed.

The reason is direct. The act of one actor sending a message to another
is not in either journal as a program statement. The sender's journal
contains the operations leading up to the send; the receiver's journal
contains the operations following its receipt. The send itself --- the
event that joins them --- has no representation in either record. Replay
can therefore reconstruct each actor in isolation, but not the causal
chain that connects them.

This consequence matters wherever event sourcing, time-travel debugging,
or post-hoc analysis is intended to span more than one actor. The model
that promises replayability delivers it within actors and not across
them.

\subsubsection{5.3 Cross-datacenter replication loses program-level
causation}\label{cross-datacenter-replication-loses-program-level-causation}

When journals are replicated across data centers --- a standard pattern
for disaster recovery and geographic distribution --- the per-actor
histories travel with them. The cross-actor causation does not, because
it was never written anywhere that replication can carry. It lived in
the dispatcher, the message broker, the runtime --- components that are
typically per-cluster or per-DC.

Recovery from a DC failure can therefore reconstruct each actor's local
state but cannot reconstruct, from the replicated material alone, the
causal chain that led to the system's pre-failure state. The chain is
reconstructed by replaying the broker's queue, by re-issuing messages
that were in flight, or by reconciliation processes that assume the
state and reconstruct backward. None of these steps is part of any
actor's program; all of them depend on infrastructure-level apparatus
that may or may not have been replicated.

\subsubsection{5.4 Debugging across actors is log
archaeology}\label{debugging-across-actors-is-log-archaeology}

A developer asked to explain why an actor produced a particular state,
when that state was caused by a chain of events crossing several actors,
cannot read a single program to find the answer. The developer assembles
the answer from logs, traces, broker dumps, and timestamps --- pieces of
evidence whose joining is itself an act of reconstruction. The
discipline that emerges is log archaeology: the careful inference of a
causal narrative from artifacts that are not narrative themselves.

The discipline is real and, in some organizations, mature. Its maturity
is proportional to the absence of programmatic causation. But its
existence is the symptom. A program whose causation is recorded as
program admits debugging by reading; a program whose causation is
dispersed across infrastructure admits debugging only by inference.

\subsubsection{5.5 Closing}\label{closing-1}

The four consequences are not separate problems with separate fixes.
They are four manifestations of a single structural absence: causation
between actors is held outside any actor's program, so any operation
that requires reading the cross-actor causal chain --- auditing,
replaying, replicating, debugging --- must reconstruct it from artifacts
that lie outside the program.

{\def\LTcaptype{none} % do not increment counter
\begin{longtable}[]{@{}
  >{\raggedright\arraybackslash}p{(\linewidth - 4\tabcolsep) * \real{0.3333}}
  >{\raggedright\arraybackslash}p{(\linewidth - 4\tabcolsep) * \real{0.3333}}
  >{\raggedright\arraybackslash}p{(\linewidth - 4\tabcolsep) * \real{0.3333}}@{}}
\toprule\noalign{}
\begin{minipage}[b]{\linewidth}\raggedright
Symptom
\end{minipage} & \begin{minipage}[b]{\linewidth}\raggedright
What is lost
\end{minipage} & \begin{minipage}[b]{\linewidth}\raggedright
What external apparatus reconstructs it
\end{minipage} \\
\midrule\noalign{}
\endhead
\bottomrule\noalign{}
\endlastfoot
Auditability & the causal chain as part of any program & distributed
tracing, correlation IDs, log aggregation \\
Replay coherence & the joint history reconstructible from per-actor
journals & custom event-sourcing layers that span actors \\
Cross-DC replication & program-level causation when journals cross DCs &
broker-level log replication, reconciliation processes \\
Debugging & the ability to read the cross-actor flow as a single program
& log archaeology --- manual correlation of timestamps and IDs \\
\end{longtable}
}

The cost of the assumption is not a list of operational difficulties to
be addressed by better tools. It is the systematic displacement of
causation from program to infrastructure. The tools that arise ---
tracing, correlation IDs, reconciliation processes, log archaeology ---
are not extensions of the program. They are evidence of what the program
does not contain.

\begin{center}\rule{0.5\linewidth}{0.5pt}\end{center}

\subsection{6. Why existing approaches do not dissolve the
assumption}\label{why-existing-approaches-do-not-dissolve-the-assumption}

The consequences of the assumption named in §3 --- auditability through
external reconstruction, replay limited to single actors, cross-DC
fragility, debugging by log archaeology --- are not unrecognized in the
actor-systems literature. Each has accumulated a corresponding
repertoire of patterns intended to address it. This section examines the
principal patterns and shows that they are responses to the
consequences, not dissolutions of the assumption that produces them.
Each pattern reconstructs cross-actor flow somewhere; none records it as
a sentence of any participating actor's program.

\subsubsection{6.1 Sagas}\label{sagas}

The saga pattern, in both its orchestrated and choreographed forms,
addresses the problem of multi-actor business transactions: an operation
that requires several actors to act in sequence, with compensating
actions if any step fails.

In the orchestrated form, a saga orchestrator holds a state machine that
represents the cross-actor flow. The orchestrator dispatches commands to
the participating actors and listens for their responses; if a step
fails, the orchestrator emits compensating commands that undo or correct
earlier steps. The cross-actor flow is therefore programmatic, but it is
the orchestrator's program --- not the program of any participating
actor. From the perspective of an actor receiving the orchestrator's
command, the message is indistinguishable from any other; the actor's
program does not know it is in a saga.

In the choreographed form, the orchestrator is dissolved into a
protocol. Each participating actor publishes events when it completes a
step; the next actor in the flow subscribes to those events and triggers
its own step. The cross-actor flow is encoded across the actors as a
distributed protocol --- but no single actor's program contains the
flow. Each actor's program contains its local response to events; the
flow itself is the implicit composition of those responses, observable
only from outside.

Both forms treat the cross-actor flow as a separate concern from any
participating actor's program. The orchestrated form externalizes it to
a coordinator; the choreographed form distributes it across event
subscriptions. Neither form makes the cross-actor send a sentence of the
sender's program. The assumption stated in §3 is preserved by
construction. The saga makes cross-actor flow programmatic only by
relocating it outside the actors whose behavior it coordinates.

\subsubsection{6.2 Distributed tracing}\label{distributed-tracing}

Distributed tracing --- OpenTelemetry, Zipkin, Jaeger, and similar
frameworks --- addresses the problem of observability across actor
boundaries. A trace context is propagated through messages by the
runtime; spans are emitted by participating components; a trace storage
backend reconstructs the cross-actor causal chain from the spans,
presenting it as a tree or timeline.

The reconstruction is real and useful. It also confirms the structural
absence it compensates for. The trace is built from artifacts that exist
outside the participating actors' programs: the trace context is
metadata threaded through messages by middleware; the spans are emitted
by instrumentation that observes the runtime; the trace storage is a
separate service that joins the spans. None of these artifacts is part
of any actor's program. They are the apparatus that makes the absence
navigable.

A trace is also lossy in a way that program is not. A trace records that
a span with given attributes occurred; a program records what was said
and why. The two are not equivalent. The trace is sufficient for
observability; it is not the program of the cross-actor flow. It exists
precisely because no such program exists.

\subsubsection{6.3 Workflow engines}\label{workflow-engines}

Workflow engines --- Temporal, Cadence, AWS Step Functions, Camunda,
Argo Workflows --- externalize the cross-actor flow into a dedicated
programming environment. The workflow is written as a separate program
in the engine's own model; the engine dispatches activities to
participating services; the engine's program records the flow as it
executes.

Workflow engines are unusual among the patterns considered here because
they do produce a programmatic record of the cross-actor flow. The flow
is a program --- an explicit, replayable, durable program. But the
program belongs to the workflow engine, not to any participating actor.
The actors are activities invoked by the workflow; their programs do not
contain the workflow.

The asymmetry matters. From the perspective of a participant, an
activity invocation is a remote procedure call from an external system;
the participant's program records that it received an invocation and
produced a result. The relationship between activities --- the flow that
connects them --- lives in the workflow engine, in a separate codebase,
in a separate execution model. Reading the participant's program does
not reveal the workflow; reading the workflow does not reveal the
participant's local program. Two complementary records that do not
compose into one. The workflow and the actors form a bipartite
description of what, in a semantically continuous system, would be a
single program.

\subsubsection{6.4 Closing}\label{closing-2}

The four patterns examined --- orchestrated sagas, choreographed sagas,
distributed tracing, and workflow engines --- are different responses to
the same consequence: cross-actor flow is not in any participant's
program, and something must compensate for that absence. The patterns
differ in where they place the compensation: in a coordinator's state
machine, in an event-subscription protocol, in a trace storage backend,
in a workflow engine's separate program.

{\def\LTcaptype{none} % do not increment counter
\begin{longtable}[]{@{}
  >{\raggedright\arraybackslash}p{(\linewidth - 4\tabcolsep) * \real{0.3333}}
  >{\raggedright\arraybackslash}p{(\linewidth - 4\tabcolsep) * \real{0.3333}}
  >{\raggedright\arraybackslash}p{(\linewidth - 4\tabcolsep) * \real{0.3333}}@{}}
\toprule\noalign{}
\begin{minipage}[b]{\linewidth}\raggedright
Pattern
\end{minipage} & \begin{minipage}[b]{\linewidth}\raggedright
Where the flow is encoded
\end{minipage} & \begin{minipage}[b]{\linewidth}\raggedright
Does it preserve the flow as program in the actors?
\end{minipage} \\
\midrule\noalign{}
\endhead
\bottomrule\noalign{}
\endlastfoot
Saga (orchestrated) & In the orchestrator's state machine & No ---
externalized to a coordinator \\
Saga (choreographed) & Across event handlers in each actor & No --- only
local responses appear \\
Distributed tracing & In the trace storage, observed by an external
system & No --- the program is observed, not extended \\
Workflow engine & In the workflow engine's separate program & No --- the
workflow is a separate artifact \\
\end{longtable}
}

The four ``No''s share a structure. In every case, cross-actor flow is
made programmatic only by placing the program somewhere other than in
the actors whose behavior constitutes the flow. The compensation always
occurs outside the participants.

These patterns are therefore not solutions to the assumption stated in
§3; they are architectural responses that accept the assumption and
build systems around it. Their sophistication, maturity, and widespread
adoption are not evidence that the assumption is correct. They are
evidence that the absence of semantic continuity is costly enough to
justify entire subsystems dedicated to reconstructing what the program
does not contain.

\begin{center}\rule{0.5\linewidth}{0.5pt}\end{center}

\subsection{7. Reformulation}\label{reformulation}

The diagnosis is complete. The assumption stated in §3 --- that
causation between actors is treated as operational rather than
programmatic --- is not entailed by the actor model (§4); produces a
structural absence of semantic continuity at every actor boundary (§5);
and survives in every existing pattern that responds to the resulting
consequences (§6). The remainder of the paper takes the diagnosis as
established and asks what would have to be the case for the assumption
to be dissolved. This section answers that question by deriving the
conditions under which semantic continuity is preserved across actors,
without violating any of the actor model's structural commitments and
without introducing orchestration.

\subsubsection{7.1 Three conditions}\label{three-conditions}

The conditions are derived directly from §4. The actor model entails
autonomy, message-based communication, and isolation (§4.1). The
construct introduced in §1 requires that the cross-actor causal chain be
recorded as part of some actor's program (semantic continuity).
Reconciling the two yields three conditions, each of which addresses one
structural commitment of the model and one element of the construct. The
conditions are not design choices; they are the minimal consequences of
reconciling the model's commitments with the construct.

\textbf{Condition C1 --- Locality of writes.} When an actor sends a
message to another actor, the sender records the act of sending in its
own journal, and only there. The receiver, on processing the message,
records the receipt in its own journal, and only there. No actor writes
to another's state. This condition preserves the model's isolation
commitment: each journal is local property; recording cross-actor
causation does not require shared state.

\textbf{Condition C2 --- Causation as program statement.} The
cross-actor send appears as a sentence in the sender's program --- a
journal entry whose content names the recipient and the message. Reading
the sender's journal reveals not only what the actor did locally but
also the act of speaking that connected its program to another actor's.
The cross-actor edge becomes part of the program's narrative rather than
an artifact reconstructed from outside it. The construct introduced in
§1 --- semantic continuity --- is realized at the sender's boundary by
this condition; symmetrically, the receiver's program records the
receipt as a sentence about who spoke to it.

\textbf{Condition C3 --- No external coordinator.} No party outside the
participating actors decides what each actor does next. Each actor
processes its own messages autonomously and decides its own response.
The journal records what happened; no orchestrator prescribes what
should happen, and no participant outside the actors is required to
interpret the record. This condition preserves the absence of
orchestration that §4.4 identified as a non-negotiable property of the
alternative.

The three conditions are independent. C1 alone preserves isolation but
does not establish continuity; C2 alone establishes continuity but,
without C1, would risk shared-state writes; C3 alone preserves autonomy
but does not address the recording question. Together, the three
conditions yield a system in which cross-actor causation is recorded as
program in each participant's local journal, dispatched through the
existing message-passing layer, and coordinated by no external party.

\subsubsection{\texorpdfstring{7.2 \emph{tell} as
primitive}{7.2 tell as primitive}}\label{tell-as-primitive}

A primitive that satisfies the three conditions by construction can be
defined.

\begin{quote}
Define \emph{tell} as a sentence in an actor's program that, in a single
act, (a) records an entry in the sender's own journal naming the
recipient and the message, (b) dispatches the message through the
existing message-passing layer, and (c) requires no coordination with
any party outside the sender and the recipient.
\end{quote}

The three components correspond to the three conditions. Component (a)
satisfies C1 --- the journal entry is written only to the sender's
journal --- and C2 --- the entry is the program statement that names the
cross-actor causation. Component (b) inherits the existing
message-passing semantics of the actor model; nothing about the dispatch
mechanism is new. Component (c) satisfies C3 --- the act involves only
the sender's local write and the dispatch; no orchestrator participates.

The primitive is conceptual. Its naming follows actor-systems
convention: in Akka and elsewhere, \emph{tell} designates a
fire-and-forget cross-actor send (§2.3). The collision with Akka's verb
is not accidental but argumentative: the verb has been part of
actor-systems vocabulary for over a decade alongside the assumption that
the verb's effect was extra-programmatic. The primitive defined here
keeps the verb and changes what the verb records. Where Akka's
\emph{tell} leaves no trace in either actor's program (§2.3), this
primitive's effect is to make the trace the program. The difference is
not in the dispatch semantics but in what the program is allowed to say
about the dispatch.

The canonical effect can be stated in one sentence: \emph{tell} turns
the journal into the boundary of causation. The boundary between actors
does not disappear; it remains the locus of message passing. What
changes is that the boundary becomes inscribed in each participating
actor's program.

\subsubsection{7.3 What changes, what does
not}\label{what-changes-what-does-not}

A reader familiar with the actor model may ask whether the conditions
above amount to a different model. They do not. Each of the actor
model's three commitments survives unchanged.

\textbf{Autonomy is preserved.} Each actor still processes its own
messages serially. \emph{tell} introduces no requirement of
synchronization with other actors; the sender does not wait for the
recipient. The conditions add a write to the sender's local journal;
they do not couple the sender's progression to the recipient's.

\textbf{Message-based communication is preserved.} Actors still
communicate by passing messages, exclusively. \emph{tell} uses the same
message-passing layer that the actor model already specifies; no shared
memory, no remote procedure calls, no synchronous read of another
actor's state. The dispatch mechanism is unchanged; only the recording
at each end is added.

\textbf{Isolation is preserved.} No actor reads or writes another's
state. The sender's journal entry is written by the sender. The
receiver's journal entry, if any, is written by the receiver. The two
records are local; they reference each other by content, not by shared
address.

What changes is what each actor's program is permitted to say about its
own activity (echoing §4.3). What does not change is what each actor can
do, observe, or share. The opacity of the cross-actor send to the
sender's program --- which §3 named as the assumption and §4
demonstrated to be contingent --- is the only thing that the conditions
remove. The actor model remains intact; only the historical
interpretation of what must remain invisible to the program is removed.

\subsubsection{7.4 The shape of the act}\label{the-shape-of-the-act}

The structural shape of \emph{tell} is rendered below. The diagram is
included not to introduce new content but to make the relations among
the components of §7.2 visible at a glance. The diagram is deliberately
mundane: nothing new is introduced except the presence of the journal
entries.

\begin{Shaded}
\begin{Highlighting}[]
\NormalTok{sequenceDiagram}
\NormalTok{    participant ProgA as Actor A\textquotesingle{}s program}
\NormalTok{    participant JA as A\textquotesingle{}s journal}
\NormalTok{    participant MPL as Message{-}passing layer}
\NormalTok{    participant ProgB as Actor B\textquotesingle{}s program}
\NormalTok{    participant JB as B\textquotesingle{}s journal}

\NormalTok{    ProgA{-}\textgreater{}\textgreater{}JA: tell B(msg) — record send entry}
\NormalTok{    ProgA{-}\textgreater{}\textgreater{}MPL: dispatch msg}
\NormalTok{    MPL{-}\textgreater{}\textgreater{}ProgB: deliver msg}
\NormalTok{    ProgB{-}\textgreater{}\textgreater{}JB: process and record receipt entry}
\end{Highlighting}
\end{Shaded}

Two journal writes (in JA and JB), one dispatch through the existing
message-passing layer, no participant outside the sender and the
recipient. The diagram reproduces the three conditions visually: writes
are local (C1), the cross-actor flow is recorded as program at both ends
(C2), and no coordinator is present (C3).

\subsubsection{7.5 Closing}\label{closing-3}

The conditions are not extensions of the actor model; they are the
conditions under which the assumption named in §3 can be removed without
altering the model's structural commitments. \emph{tell} is one
realization of the conditions; any realization that satisfies C1, C2,
and C3 would dissolve the assumption identified in §3. An actor
framework whose programs do not currently record cross-actor sends
could, in principle, be extended to satisfy the three conditions.

The reformulation is conceptual. The remainder of the paper presents the
empirical question: can the conditions be realized in a system that runs
in production? §8 answers by exhibiting an instantiation and
demonstrating it through a comparative case study.

\begin{center}\rule{0.5\linewidth}{0.5pt}\end{center}

\subsection{8. Instantiation}\label{instantiation}

\subsubsection{8.0 Origins of the
instantiation}\label{origins-of-the-instantiation}

The instantiation discussed below is provided by Puppeteer, an
actor-based framework whose journal-as-program substrate predates the
analysis presented in this paper. The framework's lineage runs from a
2005 autopersistence prototype --- domain classes that persisted
themselves through reflection, with no schema decisions in the domain
code --- through a DSL that emerged as the persistence substrate, to
event sourcing as the runtime's primary discipline. By the time the
conditions of §7 were articulated, the framework already satisfied them;
the present section reads as observation of that alignment, not as
derivation of the framework from the conditions.

\subsubsection{8.1 The case study domain}\label{the-case-study-domain}

The case study is a simplified loyalty domain modelled on a production
scenario from a deployment of the framework. A Seller actor confirms a
purchase order via a domain command. A RewardEngine actor holds a
registry of campaigns; for each purchase event, the RewardEngine
evaluates which campaigns qualify (by date and amount thresholds) and
applies the corresponding rewards to the customer.

The flow is intentionally minimal: one actor produces a domain event,
another reacts to it, and the relationship between the two needs to
cross an actor boundary. The simplicity is deliberate --- the case study
illustrates a structural property of the cross-actor mechanism, not the
complexity of any particular business workflow.

\subsubsection{\texorpdfstring{8.2 The instantiation: \emph{tell} in
Puppeteer}{8.2 The instantiation: tell in Puppeteer}}\label{the-instantiation-tell-in-puppeteer}

Puppeteer's Reaction surface exposes three planes --- \emph{Program} for
in-actor read-only effects, \emph{Metadata} for journal-metadata
changes, and \emph{Causation} for cross-actor causation. The three
planes name what the verb touches; \emph{tell} lives on the third.

The Seller's Reaction is defined as follows:

\begin{verbatim}
seller.Reactions.DefineReaction("PurchaseFunnelToRewards")
    .Job().Company().ReadForward()
    .WithSharedHydration()
    .Seek("Purchase")
        .OnMatch("[s:Seller].purchase($orderId, $date, $amount, $customer)")
    .Causation.Continue(@"
        tell RewardEngine('rewards-1')
            PurchaseConfirmed(@orderId, @date, @amount, @customer, 'STORE-42')
            id 'tid-purchase-100'
            through 'Kafka:loyalty-v1';
    ");
\end{verbatim}

The \texttt{tell} statement is a sentence in the actor's DSL. It names
the recipient
(\texttt{RewardEngine(\textquotesingle{}rewards-1\textquotesingle{})}),
the message (\texttt{PurchaseConfirmed(...)}), an envelope identifier
(\texttt{id\ \textquotesingle{}tid-purchase-100\textquotesingle{}}), and
an optional transport hint
(\texttt{through\ \textquotesingle{}Kafka:loyalty-v1\textquotesingle{}}).
When the Seller's domain command \texttt{s.purchase(...)} is executed
and Reactions are subsequently evaluated, the matching Reaction's
\texttt{.Causation.Continue(...)} body fires and the \texttt{tell}
writes one journal entry on the Seller's side.

After the bridge delivers the envelope to the RewardEngine and the
receiver acknowledges, the Seller's journal contains three entries:

\begin{verbatim}
[0]  s = Seller(); s.purchase('ord-100', 5/9/2026, 250, 'cust-42');
[1]  tell RewardEngine('rewards-1') PurchaseConfirmed(orderId, date, amount, customer, 'STORE-42')
         id 'tid-purchase-100' through 'Kafka:loyalty-v1';
[2]  tell ack 'tid-purchase-100' from RewardEngine('rewards-1');
\end{verbatim}

The three conditions of §7 are realized by these entries.

\begin{itemize}
\tightlist
\item
  \textbf{C1 (Locality of writes).} The three entries are in the
  Seller's journal only. The RewardEngine's journal records its
  receiving operation independently --- the two journals never share
  storage.
\item
  \textbf{C2 (Causation as program statement).} Entry {[}1{]} is the
  cross-actor send rendered as a DSL sentence. Reading the Seller's
  journal alone reveals that it spoke to RewardEngine, with what
  message, at what time. Entry {[}2{]} closes the round-trip with the
  receiver's acknowledgment.
\item
  \textbf{C3 (No external coordinator).} The bridge in the test is a
  transport adapter that drains envelopes and routes them to the
  receiver. It does not direct what each actor does next; it delivers
  messages produced by the actors' own programs.
\end{itemize}

The framework rejects \texttt{tell} outside
\texttt{.Causation.Continue(...)}. A direct \texttt{tell} from a
top-level command throws a runtime exception, with the error message
pointing the developer at the correct surface. The cross-actor send is
always a sentence in some Reaction's Causation body, never a
free-floating call.

\subsubsection{8.3 Three implementations side by
side}\label{three-implementations-side-by-side}

The same domain --- the Seller confirms a purchase, the RewardEngine
applies qualifying campaigns --- is exercised in three styles. The code
that distinguishes each style and the journals each produces are
reproduced below without commentary. The interpretation follows in §8.4.

\paragraph{Style 1 --- Saga
(orchestrated)}\label{style-1-saga-orchestrated}

A SagaCoordinator actor drives the workflow via direct commands to
participants:

\begin{verbatim}
saga.PerformCmd("step = 'PurchaseRequested';");
seller.PerformCmd("s = Seller(); s.purchase('ord-100', 5/9/2026, 250, 'cust-42');");

saga.PerformCmd("step = 'PurchaseConfirmed';");
rewards.PerformCmd(@"
    for (c: loyalty.Campaigns()) {
        if (c.Applies(5/9/2026, 250) == true) {
            c.Reward('ord-100', 'cust-42');
        };
    };
");

saga.PerformCmd("step = 'RewardsApplied';");
\end{verbatim}

After the run, the three journals contain:

\begin{verbatim}
SagaCoordinator's journal (3 entries):
  [0] step = 'PurchaseRequested';
  [1] step = 'PurchaseConfirmed';
  [2] step = 'RewardsApplied';

Seller's journal (1 entry):
  [0] s = Seller(); s.purchase('ord-100', 5/9/2026, 250, 'cust-42');

RewardEngine's journal (2 entries):
  [0] loyalty = RewardEngine(); loyalty.AddCampaign(...);
  [1] for (c: loyalty.Campaigns()) { ... c.Reward(...); };
\end{verbatim}

\paragraph{Style 2 --- Choreography (event-driven, no
coordinator)}\label{style-2-choreography-event-driven-no-coordinator}

An event bus mediates the cross-actor handoff:

\begin{verbatim}
bus.Subscribe(ev => {
    if (ev.StartsWith("PurchaseConfirmed:")) {
        rewards.PerformCmd(@"
            for (c: loyalty.Campaigns()) {
                if (c.Applies(5/9/2026, 250) == true) {
                    c.Reward('ord-100', 'cust-42');
                };
            };
        ");
    }
});

seller.PerformCmd("s = Seller(); s.purchase('ord-100', 5/9/2026, 250, 'cust-42');");
bus.Publish("PurchaseConfirmed:ord-100");
\end{verbatim}

After the run:

\begin{verbatim}
Seller's journal (1 entry):
  [0] s = Seller(); s.purchase('ord-100', 5/9/2026, 250, 'cust-42');

RewardEngine's journal (2 entries):
  [0] loyalty = RewardEngine(); loyalty.AddCampaign(...);
  [1] for (c: loyalty.Campaigns()) { ... c.Reward(...); };

Bus log (1 entry):
  published: PurchaseConfirmed:ord-100
\end{verbatim}

\paragraph{Style 3 --- Tell (Puppeteer)}\label{style-3-tell-puppeteer}

A Reaction on the Seller observes its own purchase and issues a
\texttt{tell} from its \texttt{.Causation.Continue(...)} body:

\begin{verbatim}
seller.Reactions.DefineReaction("PurchaseFunnelToRewards")
    .Job().Company().ReadForward()
    .WithSharedHydration()
    .Seek("Purchase")
        .OnMatch("[s:Seller].purchase($orderId, $date, $amount, $customer)")
    .Causation.Continue(@"
        tell RewardEngine('rewards-1')
            PurchaseConfirmed(@orderId, @date, @amount, @customer)
            id 'tid-comp-100';
    ");

seller.PerformCmd("s = Seller(); s.purchase('ord-100', 5/9/2026, 250, 'cust-42');");
seller.Reactions.Execute();
// bridge delivers the envelope to the RewardEngine and acks back
\end{verbatim}

After the run:

\begin{verbatim}
Seller's journal (3 entries):
  [0] s = Seller(); s.purchase('ord-100', 5/9/2026, 250, 'cust-42');
  [1] tell RewardEngine('rewards-1')
        PurchaseConfirmed(orderId, date, amount, customer)
        id 'tid-comp-100';
  [2] tell ack 'tid-comp-100' from RewardEngine('rewards-1');

RewardEngine's journal (2 entries):
  [0] loyalty = RewardEngine(); loyalty.AddCampaign(...);
  [1] for (c: loyalty.Campaigns()) { ... c.Reward(...); };
\end{verbatim}

\subsubsection{8.4 Comparative analysis}\label{comparative-analysis}

The three labs in §8.3 exercise the same logical flow but record it
differently. Three pairs of journals were exhibited.

\textbf{Saga}: the SagaCoordinator's journal contains the workflow
narrative ---
\texttt{PurchaseRequested\ →\ PurchaseConfirmed\ →\ RewardsApplied}. The
Seller's journal records its local purchase only; it does not know it is
part of a saga. The RewardEngine's journal records its local reward
only; equally unaware. Three journals, three local stories. Only the
coordinator's journal contains the joint history.

\textbf{Choreography}: no coordinator exists. The Seller's journal
records the local purchase; the publish to the bus is invisible to the
actor. The RewardEngine's journal records the local reward; the receipt
from the bus is invisible to the actor. The bus's own log records the
publish. No actor's journal contains the joint history; the bus log, an
external infrastructure artifact, is the only place the cross-actor
handoff is recorded.

\textbf{Tell}: the Seller's journal contains the purchase, the tell, and
the ack --- three entries that constitute the joint history as a
sequence of DSL sentences. The RewardEngine's journal records its local
reward, as in the other styles. Under tell, the sender's journal alone
reconstructs the cross-actor causal chain.

The three styles produce equivalent business outcomes. They differ
structurally in where the cross-actor causal chain is recorded. The
difference is not cosmetic. The saga coordinator, the event bus log, the
distributed trace, and the workflow engine all exist to compensate for
the absence identified in §3. In the tell style, the compensations have
nothing to compensate for: the program-level absence they were designed
to address is itself absent. The sender's journal already contains the
cross-actor narrative that those patterns attempt to reconstruct
elsewhere.

{\def\LTcaptype{none} % do not increment counter
\begin{longtable}[]{@{}
  >{\raggedright\arraybackslash}p{(\linewidth - 4\tabcolsep) * \real{0.3333}}
  >{\raggedright\arraybackslash}p{(\linewidth - 4\tabcolsep) * \real{0.3333}}
  >{\raggedright\arraybackslash}p{(\linewidth - 4\tabcolsep) * \real{0.3333}}@{}}
\toprule\noalign{}
\begin{minipage}[b]{\linewidth}\raggedright
Style
\end{minipage} & \begin{minipage}[b]{\linewidth}\raggedright
Joint history location
\end{minipage} & \begin{minipage}[b]{\linewidth}\raggedright
Audit path
\end{minipage} \\
\midrule\noalign{}
\endhead
\bottomrule\noalign{}
\endlastfoot
Saga (orchestrated) & In the saga coordinator's journal exclusively &
Read the coordinator's journal; participants' journals are
half-stories \\
Choreography (event-driven) & In no actor's journal; the bus log is the
only joint artifact & Merge participants' journals via correlation id,
plus the bus's log \\
Tell (program-level cross-actor primitive) & In the sender's journal as
a sequence of DSL sentences & Read the sender's journal entries
{[}tell{]} and {[}ack{]} verbatim \\
\end{longtable}
}

In the first two styles, an additional architectural element is required
to make the cross-actor flow observable as a narrative: a coordinator, a
bus log, a trace backend, or a workflow program. In the tell style, no
additional element is introduced. The narrative exists where the actor
model already records program: the journal.

\subsubsection{8.5 Property validation}\label{property-validation}

Three additional tests demonstrate that consequences claimed in §5 ---
auditability through external reconstruction, replay limited to single
actors, cross-DC fragility --- are reversed under tell, exhibiting
properties that would be inaccessible if the assumption named in §3 were
in force. These tests are intentionally chosen to mirror the
compensating patterns of §6: each test exhibits a property that, under
the assumption of §3, requires an external architectural pattern to
achieve.

\textbf{G1 --- Replay coherence (closes §5.2).} The first test stages an
in-flight tell: the envelope leaves the Seller but the bridge does not
deliver it before the test asserts. A fresh actor instance with the same
name, no shared transport, no live receiver, and no in-memory state
replays the journal alone. The replayed actor reconstructs the dedup
state for the in-flight envelope from journal entries. The joint history
exists as a program artifact and replay reaches it.

\textbf{G2 --- Cross-DC replication (closes §5.3).} The second test
replicates the Seller's journal entry-by-entry to a fresh actor in an
independent storage tier --- the analogue of moving across data centers.
The replicated actor, with no transport connectivity to the original
receiver, reconstructs the dedup state from the replicated bytes alone.
The cross-actor causal chain travels with the replication because it was
always recorded in a place replication can carry.

\textbf{G3 --- Audit query (closes §5.1).} The third test asks the audit
question --- \emph{why did this happen?} --- by reading the Seller's
journal directly. The cross-actor sentence is in entry {[}1{]}; the
acknowledgment is in entry {[}2{]}. The cause-effect chain is
reconstructed without distributed tracing, correlation IDs, or log
aggregation.

Each of these properties is a consequence of cross-actor causation being
recorded as program. Under saga, choreography, tracing, or workflow
approaches --- where it is not --- the same properties are reachable
only by consulting artifacts outside the participating actors.

\subsubsection{8.6 Closing}\label{closing-4}

The case study and the defensive tests together constitute the existence
proof. The conditions of §7 are realizable in production: a system in
which the cross-actor send is recorded as a sentence in the sender's
program, dispatched through the existing message-passing layer, and
coordinated by no external party can be built and exercised. The
instantiation in Puppeteer is one such system. Other realizations of the
conditions are possible (§7.5); the present section establishes only
that at least one is.

\begin{center}\rule{0.5\linewidth}{0.5pt}\end{center}

\subsection{9. Relation to previous work in this paper
series}\label{relation-to-previous-work-in-this-paper-series}

The journal exhibited in §8.3 --- a sequence of DSL sentences in the
sender's program, recording the cross-actor send and its acknowledgment
--- required several structural preconditions to be a viable substrate.
The journal had to be dense rather than porous: filled with operations,
not with type-erased payloads. It had to record the operations with
their parameter references intact, not with values inlined as literals.
It had to maintain a discipline that separates immediate from deferred
work, with a guardian for the boundary between them. Each precondition
has been the subject of prior structural analysis in the present series.

Paper 1 introduces \emph{porosity} --- the representational sparsity
that arises when domain state is recorded as serialized data structures
rather than as programmatic operations. Anti-porosity is the design
principle that the journal records what was said, not what was stored.
Without that property, the entries the reader saw in §8.3 ---
\texttt{tell\ RewardEngine(\textquotesingle{}rewards-1\textquotesingle{})\ PurchaseConfirmed(...)}
--- could not be programmatic at all; they would be opaque payloads.

Paper 2 introduces \emph{externalized parameters} as the structural
precondition under which compilation, caching, and dense journaling
become possible at all. Without externalized parameters, the journal
could not record what was said with parameter references intact ---
values would be inlined as literals, or the script would lose its
connection to the actor's symbol table. The tell sentence in §8.3
carries \texttt{@orderId}, \texttt{@date}, \texttt{@amount},
\texttt{@customer} as references rather than literals; this paper is
what makes that representation possible.

Paper 3 introduces the \emph{partition} between immediate and deferred
work, with Reactions as the guardian of the boundary. Paper 3 §6.8
identifies the cross-actor case as a design space open to the
alternative the present paper develops. The Reaction surface that Paper
3 establishes is the surface on which \texttt{tell} lives: the
\texttt{.Causation.Continue(...)} body the reader saw in §8.3 is where
the cross-actor send is permitted to appear. The present paper closes
the gap that Paper 3 named as honest limit.

The series, taken together, defends a single architectural property
under different framings:

\begin{quote}
\emph{Puppeteer preserves semantic continuity inside an actor. Tell
preserves semantic continuity across actors.}
\end{quote}

The first sentence is the joint contribution of Papers 1 through 3: the
structural conditions under which the journal can serve as a
program-level substrate for what an actor does. The second sentence is
the contribution of the present paper: the cross-actor extension of that
substrate that the reader saw exhibited in §8.3.

\begin{center}\rule{0.5\linewidth}{0.5pt}\end{center}

\subsection{10. Conclusion}\label{conclusion}

The actor model has, for fifty years, treated cross-actor causation as
an operational concern of the runtime rather than as a construct of the
program. The treatment was productive: actor systems became
fault-tolerant, scalable, and reliable precisely because the separation
between an actor's program and the message-passing layer was
operationally effective. Out of that productivity grew the ecosystem of
patterns analyzed in §6 and exhibited in §8.3 --- saga orchestrators,
choreography buses, distributed tracing, workflow engines. Each
compensates, in its own way, for the program-level absence the reader
saw in the journals.

The present paper observes that the absence is not entailed by the actor
model. It is a contingent design decision adopted by the canonical
sources, propagated through the lineage, and never seriously challenged
because it was never seriously challenged. The conditions under which
the absence can be removed --- locality of writes, causation as program
statement, no external coordinator --- preserve every structural
commitment of the actor model. A primitive that satisfies the three
conditions, whether named \emph{tell} or otherwise, makes the
cross-actor send a sentence in the sender's program; the apparatus that
exists to compensate for the absence has nothing left to compensate for.

The contribution of this paper is conceptual. The instantiation in
production is the existence proof; the existence proof is the journal
the reader saw in §8.3.

Forty years of convention are not forty years of necessity. The actor
model does not need to be replaced. Its commitments --- autonomy,
message-based communication, isolation --- survive the reformulation
intact. What changes is the historical interpretation of what must
remain invisible to the program. The interpretation, named explicitly in
§3, examined for contingency in §4, traced through its consequences in
§5, shown to require an entire ecosystem of compensating patterns in §6,
and contrasted with the alternative in §8.3, is the only thing the
present paper asks the field to revise.

\begin{center}\rule{0.5\linewidth}{0.5pt}\end{center}

\subsection{Appendix A. Code
references}\label{appendix-a.-code-references}

The labs cited in §8 are publicly available in the Puppeteer codebase.
The references below are to the test suite of the framework's
repository.

\subsubsection{\texorpdfstring{Cross-actor primitive surface ---
\texttt{Puppeteer/EventSourcing/Follower/}}{Cross-actor primitive surface --- Puppeteer/EventSourcing/Follower/}}\label{cross-actor-primitive-surface-puppeteereventsourcingfollower}

{\def\LTcaptype{none} % do not increment counter
\begin{longtable}[]{@{}
  >{\raggedright\arraybackslash}p{(\linewidth - 4\tabcolsep) * \real{0.3333}}
  >{\raggedright\arraybackslash}p{(\linewidth - 4\tabcolsep) * \real{0.3333}}
  >{\raggedright\arraybackslash}p{(\linewidth - 4\tabcolsep) * \real{0.3333}}@{}}
\toprule\noalign{}
\begin{minipage}[b]{\linewidth}\raggedright
File
\end{minipage} & \begin{minipage}[b]{\linewidth}\raggedright
What it shows
\end{minipage} & \begin{minipage}[b]{\linewidth}\raggedright
Cited in
\end{minipage} \\
\midrule\noalign{}
\endhead
\bottomrule\noalign{}
\endlastfoot
\texttt{Reaction.cs} & Reaction class, Action terminator dispatch,
replay-safe action execution & §7.2, §8.2 \\
\texttt{Planes.cs} & The three plane types (\texttt{ProgramPlane},
\texttt{CausationPlane}, \texttt{MetadataPlane}) and their property
accessors & §8.2 \\
\texttt{ReactionEngine.cs} & Pattern matching surface,
\texttt{.OnMatch(...)}, plane passthroughs & §8.2 \\
\end{longtable}
}

\subsubsection{\texorpdfstring{End-to-end labs ---
\texttt{UnitTestPuppeteer/}}{End-to-end labs --- UnitTestPuppeteer/}}\label{end-to-end-labs-unittestpuppeteer}

{\def\LTcaptype{none} % do not increment counter
\begin{longtable}[]{@{}
  >{\raggedright\arraybackslash}p{(\linewidth - 4\tabcolsep) * \real{0.3333}}
  >{\raggedright\arraybackslash}p{(\linewidth - 4\tabcolsep) * \real{0.3333}}
  >{\raggedright\arraybackslash}p{(\linewidth - 4\tabcolsep) * \real{0.3333}}@{}}
\toprule\noalign{}
\begin{minipage}[b]{\linewidth}\raggedright
File
\end{minipage} & \begin{minipage}[b]{\linewidth}\raggedright
What it shows
\end{minipage} & \begin{minipage}[b]{\linewidth}\raggedright
Cited in
\end{minipage} \\
\midrule\noalign{}
\endhead
\bottomrule\noalign{}
\endlastfoot
\texttt{TellLoyaltyE2ETests.cs} & Happy-path tell flow (Reaction defines
\texttt{tell} body; envelope dispatch + ack; journal verbatim asserts);
G1 replay coherence; G2 cross-DC replication; G3 audit query; negative
test for \texttt{tell} outside \texttt{.Causation.Continue(...)} & §8.2,
§8.5 \\
\texttt{CrossActorComparativeTests.cs} & Three-style side-by-side case
study: orchestrated saga, event-driven choreography, tell --- same
domain, three different journal locations of the joint history & §8.3,
§8.4 \\
\texttt{LoyaltyDomainStubs.cs} & Domain-side stubs for \texttt{Seller},
\texttt{RewardEngine}, \texttt{Campaign} --- kept minimal so the focus
remains on the cross-actor mechanism & §8.1 \\
\end{longtable}
}

\begin{center}\rule{0.5\linewidth}{0.5pt}\end{center}

\subsection{Code provenance}\label{code-provenance}

Source-code references in this paper resolve against the public
Puppeteer repository at commit
\href{https://github.com/alvaroNCubo/puppeteer/tree/2f31f9674a5de816bdf1bf9d8360ff218a02e4da}{\texttt{2f31f96}}
(2026-05-18). The snapshot is archived in Software Heritage under the
following persistent identifier:

\begin{verbatim}
swh:1:dir:10e7e6bad7eb77b6c2e406762026177f95c3ae92;
  origin=https://github.com/alvaroNCubo/puppeteer;
  anchor=swh:1:rev:2f31f9674a5de816bdf1bf9d8360ff218a02e4da
\end{verbatim}

Inline references of the form \texttt{file.cs:NN} (e.g.,
\texttt{ActorHandler.cs:38}) resolve against this snapshot. A reader can
construct a per-file SWHID by adding the qualifiers
\texttt{;path=\textless{}path\textgreater{};lines=\textless{}NN\textgreater{}}
to the directory SWHID above. Future commits to the repository may
renumber lines; the SWHID preserves the cited state independently of any
future change to the repository or its hosting.

\subsection{Acknowledgments}\label{acknowledgments}

The author used large language models (including Claude and ChatGPT) as
editorial assistants for language refinement, structural feedback, and
literature navigation. All original ideas, terminology, theoretical
constructs, and technical content presented in this work are solely the
author's.

\begin{center}\rule{0.5\linewidth}{0.5pt}\end{center}

\subsection{Appendix B. Bibliography}\label{appendix-b.-bibliography}

\emph{Online sources accessed at the moment of publication; dates filled
in at release.}

\begin{itemize}
\tightlist
\item
  Agha, G. (1986). \emph{Actors: A Model of Concurrent Computation in
  Distributed Systems.} MIT Press. ISBN 978-0-262-01092-7.
\item
  Akka documentation. \emph{Akka core: Interaction Patterns.}
  https://doc.akka.io/libraries/akka-core/current/typed/interaction-patterns.html
  (accessed 2026-05-09).
\item
  Armstrong, J. (2003). \emph{Making Reliable Distributed Systems in the
  Presence of Software Errors.} Doctoral dissertation, Royal Institute
  of Technology (KTH), Stockholm.
  https://erlang.org/download/armstrong\_thesis\_2003.pdf
\item
  Bernstein, P. A., Bykov, S., Geller, A., Kliot, G., \& Thelin, J.
  (2014). \emph{Orleans: Distributed Virtual Actors for Programmability
  and Scalability.} Microsoft Research Technical Report MSR-TR-2014-41.
\item
  Hevner, A. R., March, S. T., Park, J., \& Ram, S. (2004). \emph{Design
  science in information systems research.} MIS Quarterly, 28(1),
  75--105.
\item
  Hewitt, C., Bishop, P., \& Steiger, R. (1973). \emph{A Universal
  Modular ACTOR Formalism for Artificial Intelligence.} Proceedings of
  the 3rd International Joint Conference on Artificial Intelligence
  (IJCAI-73), pp.~235--245.
\item
  Hoare, C. A. R. (1978). \emph{Communicating Sequential Processes.}
  Communications of the ACM, 21(8), 666--677.
\item
  Microsoft Orleans documentation.
  https://learn.microsoft.com/en-us/dotnet/orleans/ (accessed
  2026-05-09).
\item
  Wensel, S. (2018). \emph{All For Reliability: Reflections on the
  Erlang Thesis.} DockYard Engineering Blog.
  https://dockyard.com/blog/2018/07/18/all-for-reliability-reflections-on-the-erlang-thesis
  (accessed 2026-05-09).
\end{itemize}

\begin{center}\rule{0.5\linewidth}{0.5pt}\end{center}

\begin{itemize}
\tightlist
\item
  Rivera, A. (2026). \emph{Anti-porous architecture: a unified design
  principle for CQRS + Actor + Event-Sourcing systems.} Paper 1 of this
  series. \url{01-anti-porosity.md}
\item
  Rivera, A. (2026). \emph{Program-value separability: the structural
  precondition for compilation, caching, and dense journaling in a DSL
  runtime.} Paper 2 of this series.
  \url{02-program-value-separability.md}
\item
  Rivera, A. (2026). \emph{Reactions and the partition: opt-in eventual
  consistency in actor-native systems.} Paper 3 of this series.
  \url{03-reactions-and-partition.md}
\end{itemize}

\end{document}
